\documentclass{jsarticle}
\begin{document}

\begin{flushleft}
睦世お姉さんへ
\end{flushleft}
\begin{flushleft}
拝啓
\end{flushleft}

\begin{center}
	\LARGE{Zeta functionについて}
\end{center}


\noindent
　  この度は、私の論文を査読してくれているかどうかはわかりませんが、
大善か火竹とは、この論文は関係なく速読でも音読でもいいので、
査読してほしいと、この度送らせてもらいました。
高等学校２年生のときに、一割る一は、三割る三は、答えが同じだが、
求めている過程の理由が違うというのと、物質量を教えてもらえてから、
数学が代替トークンとして、私が小学校の頃の算数ができなかったのが、
この商代数を教えてもらえてから、次元の単位が計算にものすごく
これは、どの数値を土台にして計算しているかが、わかり、
収率と速度、加速度、微積までもわかることができました。
その集大成が、大域的微分方程式となりました。
そのために、商代数がわかると、代数ができると、算数までもできてしまう、
その奥深さに、真逆をスポーツだけでなく、勉強までも、真逆は
すべて乗り越えることができると、この商代数が証明していました。
何回も言って恐縮ではありますが、2010年の8月に、tuplespaceによる
トポロジーのレポートをpdfで送ってから、最近まで、このアイデアが
ゼータ関数と量子群として、非可換代数方程式へと行き着くことができました。
アカシックレコードは、感情線が右手も左手もラジャループに直結していて、
三奇紋が子財と発見発明の線に旅行線を切って、両方の小指に直結しているのと、
ユウちゃんとモネとユウくんがこのラジャループと子財の源泉になっている。
家族に感謝して、どんどん、ゼータ関数のアプリケーションを作っていきます。
絶対に私を見捨てないでください。
2015年の2月28日と3月1日は、本当にすみませんでした。あんな弱気な発言
をしていたら、周りにまで、仕事と勉強のやる気と出勤の負担になるのも
よく考えていませんでした。本当にすみませんでした。

\begin{flushright}
敬具
\end{flushright}

\begin{flushright}
雅旭より
\end{flushright}


\end{document}

　
